\chapter{Elementary facts on games and regrets (Appendix A)}
\label{chap:games}

\begin{lemma}[Lemma 7]
  \label{lem:sqrt_func}
  \lean{Ito2026Adversarial.sqrt_mul_sub_mul_le}
  \leanok
  For $a, b > 0$ and $x \ge 0$, $\sqrt{a x} - b x \le \frac{a}{4b}$, with equality at
  $x = \frac{a}{4b^2}$.
\end{lemma}

\begin{proof}
  \leanok
  $\frac{a}{4b} - \sqrt{ax} + bx = \bigl(\sqrt{bx} - \frac{\sqrt a}{2\sqrt b}\bigr)^2 \ge 0$, which
  vanishes at $x = a/(4b^2)$.
\end{proof}

\begin{lemma}[Self-bounding, Lemma 8]
  \label{lem:self_bound}
  \lean{Ito2026Adversarial.le_add_sqrt_add_of_le_sqrt_add}
  \leanok
  If $x, a, b, c \ge 0$ and $x \le \sqrt{a x + b} + c$, then $x \le a + \sqrt b + 2c$.
\end{lemma}

\begin{proof}
  \leanok
  If $x \le c$ this is clear. Otherwise $(x - c)^2 \le a x + b$, so $x$ is at most the larger root
  $\frac a2 + c + \sqrt{\frac{a^2}{4} + b + ac}$ of the quadratic, and
  $\frac{a^2}{4} + b + ac \le (\frac a2 + \sqrt b + c)^2$.
\end{proof}

\begin{lemma}[Maximin mixed strategies]
  \label{lem:maximin_attained}
  \lean{Learning.ZeroSumGame.exists_forall_nashValue_le,
    Learning.ZeroSumGame.pureMaximin_le_nashValue}
  \leanok
  \uses{def:values}
  For a game with nonempty finite action sets, there is a distribution $p^\star$ with
  $u(p^\star, y) \ge \vnash$ for every $y \in \Y$, and $\vstar \le \vnash$.
\end{lemma}

\begin{proof}
  \leanok
  $g(p) = \min_y u(p, y)$ is continuous and attains its maximum on the compact simplex at some
  $p^\star$; since $\inf_q u(p, q) \le g(p) \le g(p^\star)$ for every $p$, $\vnash \le g(p^\star)
  \le u(p^\star, y)$ for every $y$. The second claim is in the library
  (\texttt{pureMaximin\_le\_nashValue}).
\end{proof}

\begin{lemma}[Relations between regrets]
  \label{lem:regret_relations}
  \lean{Learning.RepeatedGame.nashRegret_eq_psmr_add,
    Learning.RepeatedGame.nashRegret_le_externalRegret,
    Learning.RepeatedGame.IsPSNE.colGapMin_mul_le, Learning.RepeatedGame.integral_sum_rowGap_sub_le}
  \leanok
  \uses{def:regrets, lem:maximin_attained}
  For every run of a learner against an adversary:
  \begin{enumerate}
    \item $\NR_T = \PSMR_T + \Dmix T$ and $\NR_T \le \ER_T$;
    \item if $(x^\star, y^\star)$ is a PSNE and $C = \E[\#\{t < T : y_t \ne y^\star\}]$, then
      $\ER_T(x^\star) - \PSMR_T = \E[\sum_{t < T} \indic\{y_t \ne y^\star\} \DC_{y_t}] \ge C \DCmin$;
    \item for every $(x^\star, y^\star)$, if the utilities lie in $[-1, 1]$, then
      $\E[\sum_{t < T} \DR_{x_t}] - \ER_T(x^\star) \le 4 C$; if moreover $x_t$ has conditional law
      $p_t$ given the past, $\E[\DR_{x_t}] = \E[\langle \DR, p_t\rangle]$.
  \end{enumerate}
\end{lemma}

\begin{proof}
  \leanok
  \uses{lem:maximin_attained}
  (1) The first identity is linearity; for the second, with $p^\star$ of
  Lemma~\ref{lem:maximin_attained},
  $\vnash - u(x_t, y_t) \le u(p^\star, y_t) - u(x_t, y_t) = \sum_x p^\star_x (u(x, y_t) - u(x_t, y_t))$,
  so $\NR_T \le \sum_x p^\star_x \ER_T(x) \le \ER_T$. (2) $\vstar = u(x^\star, y^\star)$ for a
  PSNE, so $u(x^\star, y_t) - \vstar = \indic\{y_t \ne y^\star\} \DC_{y_t}$. (3) The summand
  $\DR_{x_t} - (u(x^\star, y_t) - u(x_t, y_t)) = u(x^\star, y^\star) - u(x^\star, y_t) + u(x_t, y_t) -
  u(x_t, y^\star)$ vanishes when $y_t = y^\star$ and is at most $4$ (no equilibrium property is
  needed).
\end{proof}

\begin{remark}
  \label{rem:no_psne_gap}
  The paper asserts that a game without PSNE has $\Dmix > 0$. This fails in general: the game with
  rows $(0, 0)$, $(1, -1)$, $(-1, 1)$ has no PSNE and $\vstar = \vnash = 0$ (row $1$ guarantees $0$,
  and a mixed strategy $p$ gets $\min(p_2 - p_3, p_3 - p_2) \le 0$). For $2 \times 2$ games it holds
  (Lemma~\ref{lem:delta_comparison}). The no-PSNE bounds of Theorems~\ref{thm:tsallis}
  and~\ref{thm:tsallis_bilinear} are therefore stated under $\Dmix > 0$, which is all their proofs
  use.
\end{remark}

\begin{lemma}[Value from equalizing strategies]
  \label{lem:value_equalizing}
  \lean{Learning.ZeroSumGame.le_nashValue_of_forall_le,
    Learning.ZeroSumGame.nashValue_le_of_forall_le,
    Learning.ZeroSumGame.nashValue_eq_of_forall_eq}
  \leanok
  \uses{def:values}
  If a distribution $p$ on $\X$ satisfies $u(p, y) \ge v$ for every $y$, then $\vnash \ge v$; if a
  distribution $q$ on $\Y$ satisfies $u(x, q) \le v$ for every $x$, then $\vnash \le v$. In
  particular, if $u(p, y) = v$ for every $y$ and $u(x, q) = v$ for every $x$, then $\vnash = v$.
\end{lemma}

\begin{proof}
  \leanok
  $u(p, q') = \sum_y q'_y u(p, y) \ge v$ for every distribution $q'$, so $\inf_{q'} u(p, q') \ge v$;
  and $\inf_{q'} u(p', q') \le u(p', q) = \sum_x p'_x u(x, q) \le v$ for every distribution $p'$.
\end{proof}

\begin{lemma}[$2 \times 2$ games without PSNE]
  \label{lem:two_by_two_value}
  \lean{Learning.ZeroSumGame.lt_and_lt_or_lt_and_lt_of_not_hasPSNE,
    Learning.ZeroSumGame.nashValue_eq_of_not_hasPSNE, Learning.ZeroSumGame.pureMaximin_fin_two}
  \leanok
  \uses{def:values, def:game, lem:value_equalizing}
  Let $u = \begin{pmatrix} a & b \\ c & d \end{pmatrix}$ have no PSNE. Then either both diagonal
  entries are larger than both off-diagonal entries ($a, d > b, c$), or the reverse
  ($b, c > a, d$); in both cases $a - b - c + d \ne 0$ and
  $\vnash = \frac{ad - bc}{a - b - c + d}$. For every $2 \times 2$ game,
  $\vstar = \max(\min(a, b), \min(c, d))$.
\end{lemma}

\begin{proof}
  \leanok
  \uses{lem:value_equalizing}
  If $a > c$, the absence of PSNE at $(1, 1)$, $(1, 2)$ and $(2, 2)$ gives in turn $b < a$,
  $b < d$ and $c < d$; if $a \le c$, the absence of PSNE at $(2, 1)$, $(2, 2)$, $(1, 2)$ and
  $(1, 1)$ gives in turn $d < c$, $d < b$, $a < b$ and $a < c$. In both cases the numerators and
  the denominator of $p^\star = \frac{(d - c, a - b)}{a - b - c + d}$ and
  $q^\star = \frac{(d - b, a - c)}{a - b - c + d}$ have the same sign, so these are distributions;
  $u(p^\star, y) = \frac{ad - bc}{a - b - c + d}$ for both columns and $u(x, q^\star)$ is the same
  value for both rows, and Lemma~\ref{lem:value_equalizing} concludes.
\end{proof}

\begin{lemma}[Lemma 10]
  \label{lem:delta_comparison}
  \lean{Ito2026Adversarial.sq_entryGap_div_four_le_mixGap}
  \leanok
  \uses{lem:two_by_two_value}
  Let $u$ be a $2 \times 2$ game with entries in $[-1, 1]$ and without PSNE. Then
  $\Dmix \ge \DM^2/4$.
\end{lemma}

\begin{proof}
  \leanok
  \uses{lem:two_by_two_value}
  Let $D = a - b - c + d$ and $V = \vnash = \frac{ad - bc}{D}$ (Lemma~\ref{lem:two_by_two_value}).
  If $a, d > b, c$, then $0 < D \le 4$, $\vstar = \max(b, c)$, and
  $V - b = \frac{(a - b)(d - b)}{D} \ge \frac{\DM^2}{4}$ and
  $V - c = \frac{(a - c)(d - c)}{D} \ge \frac{\DM^2}{4}$, each factor being the absolute value of
  an entry gap. If $b, c > a, d$, then $0 < -D \le 4$, $\vstar = \max(a, d)$,
  $V - a = \frac{(b - a)(c - a)}{-D}$ and $V - d = \frac{(b - d)(c - d)}{-D}$, and the same bound
  holds.
\end{proof}

\begin{remark}
  \label{rem:delta_comparison_statement}
  The paper assumes a unique mixed-strategy Nash equilibrium. The lemma is false for a unique
  \emph{pure} equilibrium, where $\Dmix = 0$, and the proof uses the absence of PSNE; a $2 \times 2$
  game without PSNE has a unique MSNE, so uniqueness is not assumed here.
\end{remark}
